\origin{ai}

\subsection{Forced-window carrier for the external-gauge alpha readout}
\label{sec:window6-alpha-external-gauge-d6-family-theorem-a}

\paragraph{Finite clock carrier.}
The finite clock carrier is the ten-element residue torsor
$G_{10}=\mathrm{Z}/10\mathrm{Z}$.  Its elements are residue classes $a$ modulo $10$, and its torsor operation is addition modulo $10$.  No element of $G_{10}$ is distinguished by the forced-window data alone.  A choice of an element $g\in G_{10}$ is therefore a choice of clock origin, not a theorem selecting an absolute residue name.

The external readout carrier is the product
$$
\begin{aligned}
M_{\alpha}=G_{10}\times Q_{\mathrm{space}}\times R_{\mathrm{space}}.
\end{aligned}
$$
Here $Q_{\mathrm{space}}$ is the certificate space for an externally supplied response polynomial rule, and $R_{\mathrm{space}}$ is the certificate space for an externally supplied physical readout representation.  Thus an element of $M_{\alpha}$ is a triple $(g,Q,R)$, where $g$ is a clock-origin choice, $Q$ is a response certificate, and $R$ is a representation certificate.  The carrier is finite in its clock coordinate and externally parameterized in its two certificate coordinates.

\paragraph{Displayed family.}
Let $\phi$ denote the golden ratio.  For each residue $a\in \mathrm{Z}/10\mathrm{Z}$ define
$$
\begin{aligned}
D_{6}^{(a)}=47+\phi^{-(7+a)}-\frac{1}{2}\phi^{-(17+a)}+\frac{8}{9}\phi^{-(27+a)}.
\end{aligned}
$$
The special displayed member at the zero origin is
$$
\begin{aligned}
D_{6}^{(0)}=D_{6}^{*}=47+\phi^{-7}-\frac{1}{2}\phi^{-17}+\frac{8}{9}\phi^{-27}.
\end{aligned}
$$
The ten values $D_{6}^{(a)}$ for $a=0,1,\ldots,9$ are pairwise distinct by exact symbolic simplification.  This is an algebraic statement about the ten displayed expressions, not a metrological comparison and not a physical-constant identification.

\paragraph{Two gauge actions.}
There are two distinct gauge conventions.  The residue-only gauge changes only the absolute name assigned to the coupling residue while holding the displayed exponent pattern $(7,17,27)$ fixed.  Under this convention the written expression $D_{6}^{*}$ remains the same coordinate expression.  Therefore residue-only relabeling is conservative: it treats $D_{6}^{*}$ as an externally named display, not as a forced physical value.

The exponent-shift gauge applies the clock-origin action to the exponents themselves.  A shift by $a\in \mathrm{Z}/10\mathrm{Z}$ sends the displayed exponent triple $(7,17,27)$ to $(7+a,17+a,27+a)$ and thereby produces $D_{6}^{(a)}$.  Since the ten members are pairwise distinct, this gauge action produces a ten-value family rather than a single forced readout.  The partition of the construction is therefore the ten singleton displayed values indexed by $G_{10}$, with no internally selected preferred block.

\paragraph{Local no-go statement.}
The finite-window no-go statement is the following.  The forced-window datum named here as $\mathrm{Window6}/P_{10}$ does not determine a physical value of $\alpha$, because the $G_{10}$ clock origin remains an irreducibly free gauge coordinate.  The statement uses the already certified irreducibility of the clock-origin coordinate and does not rederive that obstruction here.  Its conclusion is only that the finite forced-window structure does not select one physical alpha value from the family or from the carrier $M_{\alpha}$.

Equivalently, the construction supports the externally parameterized candidate expression
$$
\begin{aligned}
\mathrm{AlphaCandidate}_{6}^{\mathrm{ext}}=\rho_{\alpha}(\mathrm{Window6},g,Q,R),
\end{aligned}
$$
where $g\in \mathrm{Z}/10\mathrm{Z}$ and where $Q$ and $R$ are supplied by external certificates.  The map symbol $\rho_{\alpha}$ records a readout representation rule only after those certificates are present.  Without them, the forced-window data provide the gauge-indexed algebraic family and the no-go separation, but not a physical identification.

\paragraph{Certificate requirements.}
A physical readout from this carrier requires the combined certificate package
$$
\begin{aligned}
\mathrm{ClockCouplingCert}(g)\quad +\quad \mathrm{ResponseFunctorCert}\quad +\quad \mathrm{CoefficientUseCert}\quad +\quad \mathrm{PhysReadoutCert}.
\end{aligned}
$$
The first component supplies the clock-origin coupling.  The second supplies the response form used to interpret the polynomial expression.  The third justifies the use of the displayed coefficients $1$, $-1/2$, and $8/9$.  The fourth supplies the physical readout representation.  None of these four certificates is produced by the finite forced-window construction itself.

\paragraph{Witness extraction.}
The internal witness is the exact algebraic package consisting of the carrier $M_{\alpha}$, the ten-index family $\{D_{6}^{(a)}:a\in \mathrm{Z}/10\mathrm{Z}\}$, the equality for $D_{6}^{(0)}$, the residue-only invariance of the displayed expression, and the exponent-shift production of ten pairwise distinct displayed values.  The executable part of this witness uses no metrological data, no numerical value of a physical fine-structure constant, and no reverse fitting.  It is a finite symbolic certificate of the gauge family and of the no-go separation between forced-window algebra and physical readout.

\paragraph{Boundary of the construction.}
This chapter does not claim a physical fine-structure constant identification.  It does not claim that residue $7$ is forced by $\mathrm{Window6}$ or by $P_{10}$; choosing that absolute residue requires an external lexicographic convention.  It does not claim that any single $D_{6}^{(a)}$ is forward-derived as physical $\alpha$.  It does not claim that the response form $\phi^{-s}Q(\phi^{-r})$ or the coefficient use $(1,-1/2,8/9)$ is forced by the finite window.  It also does not derive a global flow law, an external dependency theorem, a metrological comparison, or an interpretation beyond the stated finite construction unless the separate certificate package supplies the missing readout data.
